\section{Supplementary Material: Model Architectures and Data Pipeline}
\label{sec:supplementary}

This section details the data preprocessing pipeline and the specific architectures of the deep learning models evaluated in the study. All models are implemented in PyTorch and were trained on a dataset of thermal video sequences.

\subsection{Data Processing Pipeline}
The input to the models consists of thermal video sequences. The preprocessing pipeline transforms raw video frames into structured tensors suitable for training.

\subsubsection{Preprocessing (Level 0 $\to$ Level 1)}
\begin{itemize}
    \item \textbf{Input}: Raw thermal video frames (MP4).
    \item \textbf{Cropping}: A Region of Interest (ROI) of size $64 \times 64$ pixels is cropped around the subject center.
    \item \textbf{Artifact Masking}: Metal reference artifacts are identified via thresholding and masked to zero to prevent model confusion.
    \item \textbf{Normalization}: Pixel temperatures are normalized using statistics computed from the training set.
\end{itemize}

\subsubsection{Input Representation}
The `SequenceHeatmapDataset` loader constructs the input tensor $X \in \mathbb{R}^{B \times T \times C \times H \times W}$ with:
\begin{itemize}
    \item $B$: Batch Size
    \item $T$: Time steps ($T=5$ frames)
    \item $H, W$: Spatial resolution ($64 \times 64$)
    \item $C$: Input Channels (Total 5):
    \begin{itemize}
        \item Channels 1-3: Thermal data (RGB duplicated or colormap mapped).
        \item Channels 4-5: Dense Optical Flow ($u, v$) computed between consecutive frames to capture heat convection dynamics.
    \end{itemize}
\end{itemize}

\subsection{Model Architectures}

\subsubsection{Standard Models (Scalar Regression)}
These models predict a vector of 4 scalar temperature values corresponding to sensor locations.
\begin{itemize}
    \item \textbf{SimpleResNet}: A frame-based ResNet-18 adapted for 5-channel input. Predicts 4 values per frame: $Y \in \mathbb{R}^{B \times 4}$.
    \item \textbf{CNN-LSTM}: Combines a ResNet encoder with an LSTM temporal aggregator. The LSTM processes the sequence of features and outputs the final temperature state: $Y \in \mathbb{R}^{B \times 4}$.
\end{itemize}

\subsubsection{Bayesian Models (Uncertainty Quantification)}
To quantify predictive uncertainty, we employ Variational Inference.
\begin{itemize}
    \item \textbf{Bayesian Neural Networks}: Standard dense layers in the regression head are replaced with `BayesLinear` layers (distributions over weights).
    \item \textbf{Output}: Returns a tuple $(\hat{Y}, D_{KL})$, where $\hat{Y}$ is the prediction and $D_{KL}$ is the Kullback-Leibler divergence regularization term.
    \item \textbf{Inference}: Uncertainty is estimated via Monte Carlo sampling (multiple forward passes).
\end{itemize}

\subsubsection{Dense Prediction and Physics-Informed Models}
\begin{itemize}
    \item \textbf{U-Net}: An encoder-decoder architecture predicting a full spatial temperature map $\hat{T}_{map} \in \mathbb{R}^{B \times 1 \times 64 \times 64}$.
    \item \textbf{Physics-Informed Loss (Bioheat)}:
    Models are trained with a composite loss function:
    $$ \mathcal{L} = \mathcal{L}_{MSE} + \lambda_{phys} \mathcal{L}_{bioheat} $$
    where $\mathcal{L}_{bioheat}$ enforces the Pennes' Bioheat equation:
    $$ \rho c \frac{\partial T}{\partial t} = \nabla \cdot (k \nabla T) - \omega_b \rho_b c_b (T - T_a) + q_m + \rho c \mathbf{v} \cdot \nabla T $$
    Here, the convection term $\mathbf{v} \cdot \nabla T$ utilizes the Optical Flow input channels.
\end{itemize}
